\begin{abstract}
High-resolution national hydrography and environmental datasets make detailed SWAT\textsuperscript{+} model construction possible across broad spatial extents, yet reproducible assembly remains difficult because routing topology, attributes, and export structure vary across analysts and study sites.
We present SWATGenX, an automated workflow that converts NHDPlus High Resolution (NHDPlus HR) and U.S.\ national ancillary layers into executable, \emph{auditable} SWAT\textsuperscript{+} project packages under an explicit, documented preprocessing rule set; here \emph{auditable} means every intermediate output is inspectable and ships with machine-readable quality-assurance sidecars.
Its central evidentiary contribution is a drainage-area fidelity audit that pinpoints \emph{where} executable model area diverges from national cumulative drainage: across 97 matched gages, SWAT\textsuperscript{+} channel areas preserve NHDPlus HR cumulative drainage at 85\% of gages within $0.5$--$2.0\times$ (and the gage-dense Peace River domain at a median ratio of $1.11$, $10$th--$90$th percentile $0.78$--$1.28$), while a stage-by-stage trace localizes a systematic $+10$--$17\%$ mainstem excess to the QSWAT\textsuperscript{+}/TauDEM channel-area computation---downstream of, and therefore exonerating, the NHDPlus HR preprocessing---a reusable, developer-relevant result about delineation-area definitions in SWAT\textsuperscript{+} (Objectives~1--2).
We further report structural and assembly-cost scaling across three size classes, split-sample calibration/verification and Morris sensitivity on two controlled gage basins, and measured simulation cost on one documented host/build (Objectives~3--5).
Reported hydrologic-skill metrics are workflow diagnostics on selected basins, not a national validation study; the contribution is reproducible, auditable model \emph{generation}, with expert review appropriate before decision use.
\end{abstract}
